\documentclass[10pt,a4paper]{amsart}
\usepackage{amsmath,amssymb,amsthm,graphicx}
\usepackage{verbatim}
\newenvironment{reference}{\comment}{\endcomment}
%\newenvironment{reference}{}{}
\newenvironment{slogan}{\comment}{\endcomment}
\newenvironment{history}{\comment}{\endcomment}

% For commutative diagrams we use Xy-pic
\usepackage[all]{xy}

% We use 2cell for 2-commutative diagrams.
\xyoption{2cell}
% Theorem environments.
%
\theoremstyle{plain}
\newtheorem{theorem}[subsection]{Theorem}
\newtheorem{proposition}[subsection]{Proposition}
\newtheorem{lemma}[subsection]{Lemma}

\theoremstyle{definition}
\newtheorem{definition}[subsection]{Definition}
\newtheorem{example}[subsection]{Example}
\newtheorem{exercise}[subsection]{Exercise}
\newtheorem{situation}[subsection]{Situation}

\theoremstyle{remark}
\newtheorem{remark}[subsection]{Remark}
\newtheorem{remarks}[subsection]{Remarks}

\numberwithin{equation}{subsection}

% Macros
%
\def\lim{\mathop{\mathrm{lim}}\nolimits}
\def\colim{\mathop{\mathrm{colim}}\nolimits}
\def\Spec{\mathop{\mathrm{Spec}}}
\def\Hom{\mathop{\mathrm{Hom}}\nolimits}
\def\Ext{\mathop{\mathrm{Ext}}\nolimits}
\def\SheafHom{\mathop{\mathcal{H}\!\mathit{om}}\nolimits}
\def\SheafExt{\mathop{\mathcal{E}\!\mathit{xt}}\nolimits}
\def\Sch{\mathit{Sch}}
\def\Mor{\mathop{\mathrm{Mor}}\nolimits}
\def\Ob{\mathop{\mathrm{Ob}}\nolimits}
\def\Sh{\mathop{\mathit{Sh}}\nolimits}
\def\NL{\mathop{N\!L}\nolimits}
\def\CH{\mathop{\mathrm{CH}}\nolimits}
\def\proetale{{pro\text{-}\acute{e}tale}}
\def\etale{{\acute{e}tale}}
\def\QCoh{\mathit{QCoh}}
\def\Ker{\mathop{\mathrm{Ker}}}
\def\Im{\mathop{\mathrm{Im}}}
\def\Coker{\mathop{\mathrm{Coker}}}
\def\Coim{\mathop{\mathrm{Coim}}}

% Boxtimes
%
\DeclareMathSymbol{\boxtimes}{\mathbin}{AMSa}{"02}

%
% Macros for moduli stacks/spaces
%
\def\QCohstack{\mathcal{QC}\!\mathit{oh}}
\def\Cohstack{\mathcal{C}\!\mathit{oh}}
\def\Spacesstack{\mathcal{S}\!\mathit{paces}}
\def\Quotfunctor{\mathrm{Quot}}
\def\Hilbfunctor{\mathrm{Hilb}}
\def\Curvesstack{\mathcal{C}\!\mathit{urves}}
\def\Polarizedstack{\mathcal{P}\!\mathit{olarized}}
\def\Complexesstack{\mathcal{C}\!\mathit{omplexes}}
% \Pic is the operator that assigns to X its picard group, usage \Pic(X)
% \Picardstack_{X/B} denotes the Picard stack of X over B
% \Picardfunctor_{X/B} denotes the Picard functor of X over B
\def\Pic{\mathop{\mathrm{Pic}}\nolimits}
\def\Picardstack{\mathcal{P}\!\mathit{ic}}
\def\Picardfunctor{\mathrm{Pic}}
\def\Deformationcategory{\mathcal{D}\!\mathit{ef}}
\title{Classification of simple field extensions: rational function field or quotient by an irreducible polynomial}
\author{The Stacks Project Authors}
\date{stacks-project:fields.tex:local-snapshot-20261012; GNU FDL 1.2 or later}
\begin{document}
\maketitle
\noindent\textit{Editorial note.} All delivered mathematical text is the author's own, copied from the pinned Stacks Project chapter file fields.tex; no statement or proof step was rewritten, reordered or removed.\par\medskip
\noindent\textit{Editorial note.} The counted proof is the author's complete argument: The source typo 'homorphism' at lines 1530-1531 of the chapter is in a different unit; within this span the author's spelling is standard.\par\medskip
\noindent\textit{Editorial note.} The label lemma-field-extension-generated-by-one-element names the counted statement, and the proof environment immediately following it in the source is delivered verbatim as its proof.\par\medskip
\noindent\textit{Editorial note.} The counted statement is carried by the author's own lemma environment in the chapter source.\par\medskip
\noindent\textit{Allowed standard prerequisites (editorial).}\begin{itemize}
\item Ring theory: the kernel of a ring morphism from k[t] is a prime ideal when the image is a domain, hence (0) or (P) for irreducible P (k[t] is a PID); universal property of quotient rings and of fraction fields
\end{itemize}\medskip
\begin{lemma}[Classification of simple extensions]
\label{lemma-field-extension-generated-by-one-element}
If a field extension $F/k$ is generated by one element, then it is
$k$-isomorphic either to the rational function field $k(t)/k$ or to one
of the extensions $k[t]/(P)$ for $P \in k[t]$ irreducible.
\end{lemma}

\begin{proof}
Let $\alpha \in F$ be such that $F = k(\alpha)$; by assumption, such an
$\alpha$ exists. There is a morphism of rings
$$
k[t] \to F
$$
sending the indeterminate $t$ to $\alpha$. The image is a domain, so the
kernel is a prime ideal. Thus, it is either $(0)$ or $(P)$ for $P \in k[t]$
irreducible.

\medskip\noindent
If the kernel is $(P)$ for $P \in k[t]$ irreducible, then the map factors
through $k[t]/(P)$, and induces a morphism of fields $k[t]/(P) \to F$. Since
the image contains $\alpha$, we see easily that the map is surjective, hence
an isomorphism. In this case, $k[t]/(P) \simeq F$.

\medskip\noindent
If the kernel is trivial, then we have an injection $k[t] \to F$.
One may thus define a morphism of the quotient field $k(t)$ into $F$; given a
quotient $R(t)/Q(t)$ with $R(t), Q(t) \in k[t]$, we map this to
$R(\alpha)/Q(\alpha)$. The hypothesis that $k[t] \to F$ is injective implies
that $Q(\alpha) \neq 0$ unless $Q$ is the zero polynomial.
The quotient field of $k[t]$ is the rational function field $k(t)$, so we get
a morphism $k(t) \to F$
whose image contains $\alpha$. It is thus surjective, hence an isomorphism.
\end{proof}

\end{document}
